\poempart{Shaping the Future}{The Power of Creation}{This part focuses on ``The Star-Absorbing Skill,'' ``The Tendon-Transforming Classic,'' and ``The Abundant Society.'' They address three questions, respectively: When intelligence becomes an external force anyone can summon, how does an individual harness it? When standard answers grow ever cheaper, how should a person train themselves? And when both capability and execution become abundant, how will organizations, division of labor, and scarcity be rewritten?

\vspace{1em}

This part is closest to where we stand, and the most \textbf{useful} of all. We discuss what happens after the rules have changed---how this artifact called artificial intelligence can better serve our lives: how to borrow its power, how to reposition our own abilities, judgment, and standing, and how to train ourselves the way we train models, becoming irreplaceable individuals who exist beyond any distribution.
}

\chapter{The Star-Absorbing Skill---When Choice Outweighs Effort}

\begin{myquote}
All conditioned phenomena are like dreams, illusions, bubbles, shadows, like dew and like lightning; thus should they be contemplated.

\hfill ---Diamond Sutra
\end{myquote}


\section{An Accountant's Puzzle: AI Arrives, Yet Her Salary Goes Up}

One day, Xiao Lin stood by the floor-to-ceiling window of her company's finance department, scrolling through the news on her phone:

{\kaiti ChatGPT passes the CPA exam, outperforming the majority of human accountants.}

She had been an accountant for ten years. Two years of professional training, eight years of hands-on practice. She could keep books, file taxes, and navigate every piece of financial software. Her boss paid her 15,000 yuan a month.

\vspace{1em}

In this second-tier city, that salary meant the market recognized her skills.

But now, an AI could finish a full day's work in three minutes---and do it more accurately. No miscalculations, no missed line items, available around the clock.

She thought she was about to lose her job.

\vspace{1em}

Three months later, her monthly salary rose to 20,000 yuan.

Not because she had mastered more advanced financial skills, but because she discovered something: the reports AI generated were flawless, yet her boss couldn't make sense of them.
What did the numbers really mean? Why had cash flow suddenly tightened? Why were accounts receivable so high? Should they adjust the procurement strategy? Should they delay the equipment investment?

\vspace{1em}

Without contextual information, AI could only produce standard answers.
This was not a technology problem---AI simply did not understand \textbf{this company's situation, here and now}.
It knew publicly available financial data from companies worldwide, knew the industry averages.

But it did not know:

{\kaiti Why had three major clients delayed payment simultaneously? (The industry had cooled---it wasn't a company problem.)}

{\kaiti Why was there an unexpected renovation expense in the warehouse? (Mandatory fire-safety compliance---unavoidable.)}

{\kaiti Why had the boss been meeting frequently with the owner of a small factory? (An acquisition was under negotiation; cash reserves needed to be prepared.)}

{\kaiti Why did the bank credit terms require positive net cash flow? (Otherwise, the credit line would be pulled next week.)}

\vspace{1em}

Xiao Lin knew all of this. Because she was on the ground every day---hearing colleagues' conversations, sensing the pressures of operations, understanding the trade-offs behind every decision.
Her job had become:

{\kaiti Have AI generate the standard financial statements (3 minutes);}

{\kaiti Cross-reference the company's actual situation, flagging anomalies that need the boss's attention (30 minutes);}

{\kaiti Write up the ``Financial Highlights'' on paper, explaining recommendations in language the boss can understand (20 minutes).}

A full day's work, now done in one hour. Yet her value had actually increased---

Because she translated data into decisions.

\vspace{1em}

Xiao Lin's story is just an ordinary example, and she was one of the lucky ones.

She was young, a quick learner, her boss was willing to give her a chance, and her role happened to be the kind AI could augment rather than fully replace. On the same floor was another colleague, Old Zhang---forty-seven years old, twenty years of bookkeeping experience. After AI arrived, the boss did not offer her a new role ``translating data into decisions.'' Instead, her name appeared on the quarterly layoff list. Her experience had not become private context; it had become sunk cost.

\vspace{1em}

Behind their stories lies a transformation that every one of us is facing or soon will face. The same transformation carries entirely different meanings for people in different circumstances. So dramatic is this upheaval that ``a once-in-a-century transformation'' is no longer a distant prophecy.
In this upheaval, some will plummet into the abyss by clinging to obsolete skills, while others will build ``one person equals a billion'' super-platforms.
Xiao Lin in our story is still at the stage where the tool helps her accelerate; what follows closely behind is---
\textbf{AI will redefine what capability means, becoming the lever that moves an era.}

\vspace{1em}

Jin Yong once wrote about a martial art that taught people how to \emph{grow stronger without grueling cultivation}.

To understand this transformation, let us return to Jin Yong's world of martial arts.

Back to Linghu Chong and his secret about ``where power truly comes from.''

\newpage
%============================================================
\section{In the AI Era, Everyone Can Become Linghu Chong, Absorbing the Power of Others}
%============================================================

When Jin Yong wrote \textit{The Smiling, Proud Wanderer}, he probably never imagined that in the AI era, everyone would effortlessly master the Star-Absorbing Skill---the supreme martial arts technique of the martial world.

The classical way of cultivation in martial arts fiction was \textbf{grueling practice}: Guo Jing, dull-witted by nature, spent decades through bitter winters and scorching summers drilling the Eighteen Dragon-Subduing Palms, one move at a time, into his very body. This kind of mastery was physiological---muscle memory, neural reflex. It could not be taken away; it lived and died with the flesh.

\vspace{1em}

The cultivation of ability followed a simple logic: \textbf{internalization}.

It was not enough to know; the body and brain had to be rewritten, again and again.

Every time Guo Jing threw the Eighteen Dragon-Subduing Palms, it left a trace in his mind. Those traces were not in the secret manual---they were in the pain of being corrected by reality, time after time. Hong Qigong could demonstrate for Guo Jing, but the final leap---from knowing to doing---had to be completed by his own body.

\vspace{1em}

It could not be transmitted. It could only be re-experienced, re-lived.

Knowing is easy; doing is hard. Most of us cannot learn to ride a bicycle without falling a few times.

This was the age of \textbf{mastery through the flesh}.

\vspace{1em}

In this age, the acquisition of ability followed a simple formula:

\begin{equation*}
\text{能力} = f(\text{天赋}, \text{时间}, \text{方法})
\end{equation*}

Here, time is conserved---it cannot be compressed. To master a complex ability, one must invest time. The ``10,000-hour rule'' is not a rigorous scientific law, but it reminds us of at least this much: no tool, however powerful, can entirely bypass long-term practice and real-world feedback.

\begin{importantbox}
    Yet the conservation of time is both fair and unfair to everyone.
\end{importantbox}

\textbf{Fair in this sense:} everyone's time is equally long.

What you put in is what you get out; our upper bound is the total time we can invest.

No matter how gifted, a person's day still has only 24 hours.

\vspace{1em}

\textbf{Unfair in this sense:} starting points are never the same.

Some are born with extraordinary gifts, able to memorize at a glance; others are of average talent, requiring a hundred rounds of tempering to achieve mastery.

Some have the guidance of great teachers, gaining direct access to true knowledge; others grope in the dark, hitting walls for half a lifetime.

Some are born into families of scholars, reading the \textit{Analects} as toddlers; others are trapped in impoverished villages, still herding cattle at age ten.

\vspace{1em}

In the age of mastery through the flesh, talent and birth were invisible accelerators.

Different starting lines lead to different finish lines.

Ability is our property; cognition is our moat.

They define who we are and determine our position in the world.

\textbf{Lose them, and we lose the self that was forged through the flesh.}

\vspace{1em}

Linghu Chong took a different path: \textbf{directly absorbing others' inner power, converting it to his own in an instant.} No need to understand how that power was cultivated; no need to repeat the predecessors' training. He simply needed to---\textbf{call upon it}.

\vspace{1em}

These are two fundamentally different modes of cultivation:

Mastery through the flesh says: I am what I \textbf{possess}.

The Star-Absorbing Skill says: I am what I can \textbf{invoke}.

\vspace{1em}

Jin Yong planted a warning: the absorbed inner forces would clash within the body, causing internal energy to go haywire. Linghu Chong was nearly torn apart by these powers ``that did not belong to him.'' This is precisely the truth of the Star-Absorbing Skill: invoking another's ability is not the same as possessing the core to wield it.

\vspace{1em}

\textbf{We will ultimately be consumed by the very power we borrowed.}

The AI era is pushing all of us into Linghu Chong's predicament.


\begin{thinkbox}[Reflection]
If ability no longer requires grueling effort to acquire, then what is the meaning of ``hard work''?

If everyone can summon the same power,

will the gap between people vanish, or will it widen in a different form?
\end{thinkbox}


\newpage
%============================================================
\section{Two Evolutions of the Star-Absorbing Skill: Internalization > Encapsulation > Diffusion}
%============================================================

Ability is the foundation on which we make our place in the world. Zoom out far enough, and we see that the very definition of ability is being rewritten---stripped away from the individual, step by step.

I summarize this process as two evolutions: ability was first separated from the body and encapsulated in tools, then diffused into the air itself. Each evolution brought a reassessment of individual value, and we are now in the midst of the second.
Let us see how ability went from ``private property'' to ``public infrastructure.''

\subsection{Legendary Weapons: Ability Encapsulated in Tools}

Humans use tools---this is hardly news, and not even uniquely human.

The stone axe extended the arm; the bow extended throwing range; the telescope magnified the eye's reach.

In the age of mastery through the flesh, tools were amplifiers: they magnified abilities one already possessed.

The age of legendary weapons changed the rules:

\begin{importantbox}
    Ability no longer had to reside in the body; its nature shifted from ``amplification'' to ``encapsulation.''
\end{importantbox}

Once, the martial world was decided by fists and inner power---plucking a flower or a leaf could wound a man, but it demanded decades of bitter training. Now, \textbf{with the ``Peacock Plume,''} even a boy with no inner power at all could pull the trigger, and in a flash of radiant light, slay a veteran master famed for decades.

\vspace{1em}

Killing power was encapsulated in precision gears and mechanisms, no longer dependent on one's inner energy reserves. Once encapsulated, ability could be bought and sold, transmitted, or seized.
We have the magical weapon; others do not---that is the advantage of a generational gap.

\vspace{1em}

Of course, the ``Peacock Plume'' was formidable, but it had its own methods of operation: one had to know how to activate the weapon, how to aim at the critical point, how to absorb the recoil when the mechanism fired. Without knowing how to commune with the weapon's spirit, without knowing how to channel its power, this legendary weapon in hand might be less useful than a common fire poker.
We know the incantation; others do not---and so the gap opens.

\subsection{Speak and It Shall Be Done: Ability Everywhere, Summoned by Thought}

The arrival of AI brings yet another transformation of ability. This time, ability neither needs to ``reside in the body'' nor to ``be housed in a tool''---it becomes a force that can be summoned at any moment.

Imagine the spiritual energy of cultivation novels---omnipresent, arriving at the mere thought of it.

No need to ``possess'' the energy, because it is already in the air.

No need to ``cultivate'' the energy, because heaven and earth have already refined it.

\textbf{One needs only to know how to breathe it in,} and like the Star-Absorbing Skill, one absorbs ability itself.

\vspace{1em}

This is the form ability takes in the AI era.

Like \textbf{food delivery}: you don't need to know how to cook---just how to place an order.

Specify the flavor, budget, dietary restrictions, delivery time, and so on.

\vspace{0.5em}

Like \textbf{navigation}: you don't need to memorize maps---just enter the destination.

Take the fastest route, avoid the highway, pick someone up along the way, and so on.

\vspace{0.5em}

Like a \textbf{personal errand service}: you don't need to visit every office window yourself---just clearly describe what you need.

What documents are required, which process to follow, what pitfalls to avoid, and so on.

\vspace{1em}

Ability has become a force of heaven and earth, omnipresent.

Simply \textbf{describe it}---speak, and it shall be \textbf{done}.

\begin{thinkbox}[Reflection]
In an age where speaking makes it so, if everyone can summon the same cosmic force,

what does that mean for each of us?

\vspace{1em}

If the ability to ask and describe a problem matters more than the ability to solve it,

might the humanities redefine what ``expertise'' itself means?
\end{thinkbox}

\newpage

\subsection{A New Starting Line: The Floor Has Been Raised, and the Ceiling Even Higher}

For every individual, the starting point jumps from 0 straight to 80. The Star-Absorbing Skill has made time far less rigidly conserved than it once was. As long as we are willing, many abilities that once required years of training just to reach the threshold can now yield a usable version in remarkably short order. Someone with zero background can dictate to AI and produce a polished business email; someone who never studied painting can use AI to quickly generate a passable image.

\vspace{1em}

Let us understand this change mathematically. As shown in Figure~\ref{fig:ability_curve}, the curve on the left is the original distribution of ability. The change introduced by AI manifests mathematically as an extreme form of ``distributional compression and polarization.''

\begin{figure}[h]
\centering
\begin{tikzpicture}
% 坐标轴
\draw[->] (0,0) -- (10,0) node[right] {\ifdefined\ENGLISH Ability\else 能力\fi};
\draw[->] (0,0) -- (0,5) node[above] {\ifdefined\ENGLISH Population\else 人数\fi};

% 原始分布（虚线）
\draw[dashed, gray, thick] plot[smooth, domain=0:8] (\x, {3*exp(-(\x-4)^2/2.5)});
\node[gray] at (4, 3.8) {\small \ifdefined\ENGLISH Original distribution\else 原始分布\fi};

% AI后的分布（实线）
\draw[thick, darkCyan] plot[smooth, domain=4.5:9.5] (\x, {3.5*exp(-(\x-6.4)^2/0.6) + 0.5/(\x-3.8)});
\node[darkCyan] at (6.4, 4.5) {\small \ifdefined\ENGLISH Post-AI distribution\else AI引入后分布\fi};

% 分数标注
\draw (0, -0.1) -- (0, 0.1);
\draw (4, -0.1) -- (4, 0.1);
\draw (6.4, -0.1) -- (6.4, 0.1);
\draw (8.5, -0.1) -- (8.5, 0.1);
\node at (0, -0.5) {\small 0};
\node at (4, -0.5) {\small 50};
\node at (6.4, -0.5) {\small 80};
\node at (8.5, -0.5) {\small 100+};

% 关键标注
\draw[->, cyan, thick] (1.5, 1) -- (5.8, 2.5);
\node[cyan] at (1.8, 1.8) {\small \ifdefined\ENGLISH Floor raised\else 地板抬高\fi};

\draw[->, cyan, thick] (7.6, 0.7) -- (9.7, 0.25);
\node[cyan] at (8.8, 1.2) {\small \ifdefined\ENGLISH Super-individuals\else 超级个体\fi};
\end{tikzpicture}
\caption{\ifdefined\ENGLISH AI reshapes the ability distribution: compressed middle, elongated tail\else AI导致的能力分布重塑：中间压缩，尾部拉长\fi}
\label{fig:ability_curve}
\end{figure}



\textbf{The floor has been raised:} AI acts as a powerful equalizer, rapidly bringing a large number of tasks---once impossible, too slow, or too inconsistent---up to a passing grade. In many highly standardized jobs, the gaps between people that once arose from differences in proficiency have been visibly compressed.

\textbf{The ceiling has been pulled even higher:} For ``super-individuals'' who can wield AI, it is not a floor but a lever of infinite magnification. While the masses crowd the plateau at 80, super-individuals have already broken through the old ceiling of 100, racing toward a height of 1,000.

\begin{importantbox}

AI has shortened the distance between mediocrity and excellence, yet widened the chasm between excellence and greatness.

The floor has been raised very high; the ceiling has been pulled even higher.
\end{importantbox}

The story of super-individuals is still distant for most of us. The reality we must first confront is this: AI has eliminated the incapacity born of skill deficiency. Its direct consequence is what I call \textbf{ability inflation}---ability, like currency, undergoes inflation.

\vspace{1em}

In a world of monetary inflation, savings lose value. In a world of ability inflation, ``competent'' loses value.

\textbf{When everyone stands at the 80-point starting line, 80 becomes the new zero.} This 80 points of ability arrived so suddenly that most people were utterly unprepared to wield it.

\vspace{1em}

One cannot help but think of Li Yuanba, the mightiest warrior of the Sui-Tang legends.

He wielded a pair of invincible golden hammers; no mortal could harm him.

Yet this young hero ultimately died by the very hammers that had made him invincible across the land---when he hurled them at the thundering sky, the force he had amplified with his own hands returned with the falling hammers and ended his life with unerring precision.

\vspace{1em}

Our arms have been extended by AI; our strength has been amplified by algorithms.

Without a matching reverence for the underlying laws and a clear sense of direction, the seemingly indestructible lever of AI in our hands will, at the next blind swing, become a boomerang aimed squarely at ourselves.

\vspace{1em}

This is precisely the change the Star-Absorbing Skill brings to our lives:

Ability is no longer scarce---direction will be.

Execution is no longer the bottleneck---judgment is.

\vspace{1em}

In the AI era, every one of us has ``absorbed'' Li Yuanba-level divine power from the models.

Let us first peer into the alchemical furnace of AI and see what this ``golden hammer'' really is, in technical terms.

\newpage
%============================================================
\section{A Technical Explanation of Why Choice Outweighs Effort}
%============================================================

An LLM is, at its core, a probability sampling machine.

When we pose a question, it generates not an answer but a distribution.

This leads to two consequences:

For any given question, there are infinitely many ``reasonable'' answers---they differ only in probability;

The model itself makes no choices---it merely tells us how ``human-like'' each answer is.

\vspace{1em}

``Human-like,'' in technical terms, means: \textbf{how frequently this expression appeared in the training data}.

\vspace{1em}

Think of it as a pachinko machine:

Drop in a question (the prompt), and it clatters out countless marbles (possible responses).

The distribution of these marbles is fixed, forming something like a bell curve:

Most land in the middle---correct, appropriate, error-free;

A few land at the extremes---either strikingly brilliant or wildly off the mark.

\vspace{1em}

The model tends to produce the peak-of-the-curve, most-intuitive-to-the-masses \textbf{statistical average}.

But the ``best answer'' we seek often lies at the edge of the long tail.

The machine's job is to reduce the cost of generating options to zero. But ``picking the option that best fits the need from among those options''---that is something the machine cannot readily do for us.
Because it requires specific context, distinctive taste, or professional judgment---information that lives not in the model's parameters but in our own minds.

\vspace{1em}

This is the new meaning of ``choice outweighs effort'': the bottleneck of productivity has shifted from ``the toil of production'' to ``the discernment of taste.''

\begin{thinkbox}[Reflection]
    An LLM is ``the average of the world,''

    but why does everyone feel that their ability has been enhanced?
\end{thinkbox}

\newpage
\subsection{``Average'' Does Not Mean ``Mediocre''---It Means ``The Standard Answer''}

A common misconception is this: since the model outputs ``average'' quality, roughly half of all people should find it inferior to themselves. The reality is the exact opposite.

Almost everyone feels their ability has been enhanced. Why?

\vspace{1em}

The reason is that we misunderstand what ``average'' means.

The ``average'' here does not mean the ability is at a middling level. Rather, it means the model consistently gives \textbf{the response closest to the standard answer}. More precisely, given a question $x$, the model generates the answer $y$ that is \emph{most common, most endorsed, and least error-prone} in the training data.

This is not a random sample of some arbitrary person's level. It is an aggregation of collective wisdom, yielding a relatively reliable reference answer.

\vspace{1em}

This ``standard answer'' quality manifests as follows:

\textbf{Writing poetry}: the meter is correct, the imagery classical, the emotional arc follows tradition---it produces ``the kind of good poem you'd find in a textbook,'' but struggles to deliver that \textbf{rule-breaking} sense of estrangement;

\textbf{Writing code}: the structure is clean, edge cases are covered, the style is consistent---it gives ``the implementation a professor would mark full score,'' but it may not be the kind of clever hack only a true geek would know;

\textbf{Presenting a viewpoint}: the argument is rigorous, the perspectives comprehensive, the phrasing airtight---it offers ``a consensus that offends no one,'' but in doing so loses the sharp insight and \textbf{paradigm-shattering perspective}.

\vspace{1em}

In other words, what an LLM produces is not some random passerby's level of work. It is \textbf{the distilled essence of thousands upon thousands of high-quality answers}. This explains why it helps the vast majority of people---because it represents ``the most reliable portion'' of humanity's knowledge base.

\subsection{How the ``Standard Answer'' Is Made, and What It Costs}

In theory, an LLM can generate any possible text: it could write \textit{Hamlet}; it could also output gibberish. It could produce Einstein's papers; it could also churn out pulp fiction. Like Borges's Library of Babel---theoretically containing every possible book, yet in practice the overwhelming majority are meaningless combinations.

To make the model useful, the training process applies \textbf{layer upon layer of selection}:

\textbf{Training data curation}: what enters the model is not all available text, but carefully selected high-quality content---academic papers, professional documents, quality web pages, and so on. This step alone dramatically reduces the probability of the model producing ``garbage text'';

\textbf{Post-training alignment}: through human feedback, preference labeling, and similar methods, the model's ``taste'' is further tuned---making it favor outputs that ``humans consider good'' and suppress those ``humans consider bad.''

The entire process can be understood through this pseudo-formula:

\begin{equation*}
p_{\text{raw}}(y|x) \xrightarrow{\text{剪枝}} p_{\text{clean}}(y|x) \xrightarrow{\text{对齐}} p_{\text{aligned}}(y|x)
\end{equation*}

Every round of selection does the same thing:

\textbf{Raise the probability of ``good answers,'' lower the probability of ``bad answers.''}

\vspace{1em}

What does this selection cost?

\textbf{The upside}: the model becomes reliable, safe, and ``obedient.'' Ask any routine question and it delivers a stable 80-point answer---never too bad, never too bold.

\textbf{The downside}: the model becomes conservative, rarely surprising. Those outputs that ``deviate from the mainstream'' are systematically suppressed---including the genius-level, the iconoclastic, the paradigm-breaking possibilities.

In other words, we traded tail risk for central stability. Among those answers that 99.9\% of people would never choose, there may lurk the 0.1\% that constitute works of epoch-making brilliance.

\begin{thinkbox}[Reflection]
When ability becomes readily available, we will enter the ``illusions'' of infinite choice.

How do we make better choices?
\end{thinkbox}

To learn to choose well, we must first see through the illusions that the Star-Absorbing Skill creates.

They do not begin with pain---they often arrive in the name of empowerment;

They rely not on coercion, but on the path of least resistance;

They do not wound us---they leave us weightless.

These illusions differ on the surface but share a common root.

\newpage
% ============================================================
\section{Three Illusions of the Star-Absorbing Skill}
% ============================================================

% ------------------------------------------------------------
% \subsection{幻境一：看尽长安，乱花迷人眼（选择太多你不会认真选）}
\subsection{Illusion One: All the Flowers of Chang'an Dazzle the Eye}

% ------------------------------------------------------------

The most immediate impact of this illusion on each of us is: when AI can generate heaps of proposals and countless possibilities in one breath, how do we choose?

It is like standing in a room with many doors, each potentially concealing treasure. We spend all our time pacing from door to door, peeking at this one, peering at that one, and in the end we open none. Meanwhile, someone else picks a door at random, walks in, finds it wrong, walks out, and eventually explores several rooms---while we are still standing in place.

\vspace{1em}

The direct cause of being dazzled is \textbf{the absence of an internal standard of judgment}.

If you know what you want, 100 options and 10 options are no different---you quickly eliminate 90\% and focus on the remaining 10\%. But if you don't know what you want---even 2 options will leave you agonizing.

\vspace{1em}

Yet having a standard alone is not enough. Even when we know what we want, a subtler poison erodes our ability to act---what I call the \textbf{illusion of reversibility}: because AI's output can be regenerated infinitely, people mistake choices themselves for being infinitely repeatable. AI makes us feel that everything can be redone---chose wrong? Generate another. Not satisfied? Have it revise. This version won't work? Try the next one. When we believe every choice can be undone and remade, we will never choose seriously.

\vspace{1em}

Yet the things in life that truly matter are precisely the irreversible ones.

Choosing this person means not choosing that one; taking this road means the time will never come back;

Words spoken are water spilled; missed opportunities do not wait around.

The illusion of reversibility drains all weight from our choices, leaving us forever gliding on a mindset of ``this will do for now, I can always change it later''---never truly committing to anything.

\begin{thinkbox}[Reflection]
Think back to a recent time you spent a long while making a decision.
Was it because there were too many options, because you lacked clear criteria, or because you left yourself a way out?
\end{thinkbox}

\newpage
% ------------------------------------------------------------
% \subsection{幻境二：纸上得来，浮光终究浅（AI会让你失去思考能力）}
\subsection{Illusion Two: What Comes from the Page Remains Only a Shimmer}
% ------------------------------------------------------------
The impact of this illusion is: \textbf{our brains are gradually losing the ability to think deeply.}

A programmer who used AI to write code for three months found that without AI, he could no longer recall even basic syntax---let alone the complex code logic.

\vspace{1em}

Is this a problem? \textbf{On the surface, no.}

Today we write code in Python and JavaScript; no one uses assembly language---and what difference does it make? The essence of technological progress is freeing people from having to master low-level details.
Vibe coding\footnote{Python, JavaScript, and assembly are traditional programming languages; Vibe coding, first proposed by Andrej Karpathy, allows programming through natural language and may become the mainstream programming paradigm of the future.} will be the Python of even more people in the future.
But this is only the surface. The real danger lies in the regression of cognitive ability.
It is not that ``we can no longer write code''---that, in itself, does not matter.

\vspace{1em}

We used to think we learned to code in order to write code, and learned to write in order to produce articles. But in reality, the process of learning these skills is the process by which the brain constructs its \textbf{cognitive frameworks}---it is not about mastering a specific operation, but about training a way of thinking.

\vspace{1em}

When we use words to deconstruct social phenomena, we are training the brain's ability to organize chaos.

When we use code to solve mathematical problems, we are training the brain's ability to decompose complexity.

The true product of this training is not code or essays---it is the brain's capacity to process problems.

\vspace{1em}

Long-term reliance on AI for answers switches the brain from ``generative mode'' to ``selection mode.'' We no longer construct logic; we merely perform low-level ``posterior screening'' among AI-provided options. When we no longer personally walk through the reasoning from zero to one, we become passengers sitting in the driver's seat who have never once gripped the steering wheel---believing we are driving, when in fact we are merely being transported.

\begin{thinkbox}[Reflection]
Think of the last time you ``felt you understood'' something, yet could not derive it from scratch.

Was it because you lacked time, or because you never walked that chain of reasoning?

If the answer is the latter, which piece of training would you be willing to pick back up?
\end{thinkbox}


% ------------------------------------------------------------
% \subsection{幻境三：逐鹿者不顾兔，却溺于流沙（习惯了AI的即时反馈而忽略真正有重要的事）}

\subsection{Illusion Three: The Deer-Hunter Ignores the Hare, Yet Drowns in Quicksand}

% ------------------------------------------------------------

The third illusion is the most insidious: becoming so accustomed to AI's instant feedback that we neglect what truly matters. We feel ourselves growing ever stronger, yet have no idea where to direct that strength. It is like someone who spent twenty years perfecting their running technique, only to wake up one day and find that teleportation has been invented. Standing before the portal, they realize: they never once thought about where they wanted to go---they had only ever practiced running faster.

% \textbf{更危险的是，当内心空了时，别的东西会趁虚而入。}
\vspace{1em}

AI excels at giving us short-term feedback:
smoother copy, more eye-catching headlines, metrics that climb more easily.
Click-through rates, likes, completion rates---anything that can be measured immediately, AI can help us optimize to perfection.

\vspace{1em}

We have already learned about overfitting in machine learning: a model that learns the training data too well loses its ability to generalize---it can only handle problems it has seen before, and the slightest variation stumps it.

\vspace{1em}

\textbf{Humans can overfit too.} When we treat likes and instant rewards as our reward function, we are training ourselves on noise---becoming ever better at winning noise, yet ever worse at winning at life.

\vspace{1em}

What is more frightening is that this overfitting is self-reinforcing: short-term feedback arrives fast, dopamine flows on time, and the brain increasingly favors this type of reward; long-term value feedback is slow and uncertain, so the brain increasingly avoids that kind of investment. Over time, the things that truly matter---the things that take ten years to show results, that cannot be quantified, that no one will ``like''---are permanently excluded from one's life.

\vspace{1em}

Bill Gates once said that people always overestimate short-term change and underestimate long-term change.

In the AI era, this observation needs an update: \textbf{people will excessively chase measurable short-term feedback, while permanently avoiding value that cannot be quantified.}

\begin{thinkbox}[Reflection]
If tomorrow all your work were done by AI and you only had to decide ``what to do,''

would you still know what you want to do?

If not, what would you use to pretend you do?
\end{thinkbox}

\newpage
% ============================================================
\section{Conclusion: The Price the Star-Absorbing Skill Exacts}
% ============================================================

These three illusions are the price the Star-Absorbing Skill so readily exacts: internal energies in turmoil.

Without your own standard of judgment, you cannot choose among the options.

Without your own framework for thinking, you can only lean on external logic.

Without your own direction in life, even the mightiest execution is just spinning in place.

\vspace{1em}

Some may ask: where does this ``inner operating system'' come from?

Our values, our criteria for judgment, our life's direction---were these not also implanted from the outside?

Shaped by parents, shaped by education, shaped by culture, shaped by society.

If everything internal was once external, then why should AI's influence be more dangerous than a parent's?

\vspace{1em}

This is a good question. My answer is:

\textbf{The difference is not between internal and external. The difference is whether it has been digested.}

\vspace{1em}

The shaping done by parents, education, and society all passed through the body's resistance and rumination.

We consented, and we rebelled.

We accepted, and we questioned.

We followed along, and we paid the price.

Through this long process of digestion, what came from outside truly became ``our own.''

\vspace{1em}

The danger of AI is that \textbf{it is too frictionless.}

It requires no digestion, triggers no resistance, leaves no room for questioning.

It glides straight into our thinking and becomes the default option.

Before we can even ask ``is this right?'' we have already hit ``send.''

\vspace{1em}

So how do we train this ``inner operating system'' for ourselves?

My answer lies in an unexpected place: train yourself the way we train AI models.

That is what Chapter Nine, \textit{The Tendon-Transforming Classic}, is about.
